% Title option 1:
% Persistence-Level Null Results for Weak Substrate-Law Reinforcement on a Kernel-Induced Cochain Operator Architecture
% Title option 2:
% A Negative Survey of Weak Classical Interaction Laws in a Kernel-Induced Cochain Packet Architecture
% Title option 3:
% Failure of Weak Substrate Reinforcement to Produce Persistence in a Discrete Cochain Packet Architecture

\documentclass[11pt]{article}

\usepackage{amsmath,amssymb,graphicx}
\usepackage[margin=1in]{geometry}
\usepackage{booktabs}
\usepackage{array}

\newcommand{\placeholderbox}[1]{%
  \fbox{\rule{0pt}{#1}\rule{0.92\linewidth}{0pt}}%
}

\title{Persistence-Level Null Results for Weak Substrate-Law Reinforcement on a Kernel-Induced Cochain Operator Architecture}
\author{Tomislav Rupi\'c \\ Independent Researcher}
\date{March 2026}

\begin{document}

\maketitle

\begin{abstract}
This paper skeleton documents the completed Phase II negative survey in the Harmonic Address Operating System (HAOS) / Interaction Invariance Physics (IIP) pre-geometry program. The underlying kernel-induced cochain operator architecture is held fixed while a broad family of weak substrate-law, gated-interaction, and delayed-memory modifications is tested against persistence-first observables. The central question is not whether bounded deformation can be produced, but whether any such deformation changes regime at the level of composite lifetime, binding persistence, coarse basin retention, corridor dwell, or stabilized ordering class. Across the Stage 19--22 program the tested mechanisms are numerically active, bounded, and non-null at the substrate or interaction layer, yet no persistence-level regime change is observed. The resulting claim is architectural and negative: weak substrate-law reinforcement, delayed memory, and selective interaction gating are insufficient, at the tested couplings and within the present model class, to generate persistent composites or anchoring. This negative survey therefore defines a structural boundary on a broad nearby weak-coupling mechanism family and motivates subsequent work on new operator sectors or new dynamical degrees of freedom rather than continued coefficient tuning.
\end{abstract}

\section{Introduction}

The Phase II program is best understood as a persistence-first falsification survey. The relevant discriminant is not whether a modified law produces visible deformation, local steering, or bounded activity. The discriminant is whether it changes the persistence ladder. In the present architecture, persistence means a sustained change in the observables that matter for composite survival: composite lifetime, binding persistence, coarse basin retention, corridor dwell or locking, and the appearance of a new stabilized ordering class. If those observables remain flat, then a law may be active without being regime-changing.

This framing matters because the Harmonic Address Operating System (HAOS) / Interaction Invariance Physics (IIP) program had already mapped a strong frozen-architecture boundary before Phase II began. The frozen architecture could organize, localize, and support repeatable collision morphologies, but it did not produce binding. Phase II therefore did not begin as a search for arbitrary novelty. It began as a constrained question: can nearby weak substrate-law or interaction-law modifications rescue persistence without changing the packet ontology or the operator class?

The paper is intentionally written as a boundary statement rather than a positive-mechanism paper. Its scientific value lies in documenting that several physically suggestive weak-law families were tested under a common metric discipline and remained persistence-null. The emphasis is therefore on scope, architecture, and falsification rather than on speculative interpretation; compare, for example, the role of discrete exterior calculus and related lattice operator formalisms in setting the architectural baseline \cite{dec_placeholder}, and the role of weak nonlinear or effective-medium analogues in motivating the tested mechanism classes \cite{nonlinear_placeholder,effective_medium_placeholder}.

\section{Architecture and Observable Layer}

The architectural baseline is a kernel-induced cochain packet dynamics with fixed support conventions, fixed representative packet families, and a fixed promotion ladder. In Phase II the architecture is not rewritten from first principles at each stage. Instead, the same representative encounter classes are reused so that the effect of each law-class modification can be evaluated against a stable comparison background. This is what makes a negative survey meaningful rather than anecdotal.

At the observable layer, the program is persistence-first. The paper should define the primary observables in words before any stage-by-stage discussion:
\begin{itemize}
\item composite lifetime: the duration over which a candidate encounter remains locally assembled,
\item binding persistence: the fraction of the observed window spent in a close or retained configuration,
\item coarse basin persistence: the survival of occupancy in a mesoscopic retention basin,
\item corridor dwell or locking: sustained transport-channel occupation or directional trapping,
\item stabilized ordering class: emergence of a new long-lived structural regime.
\end{itemize}

The promotion logic should also be stated explicitly. Visual or secondary diagnostics do not earn follow-up on their own. A branch is promoted only if persistence improves at the primary observable layer. This matters because several Phase II branches show bounded and measurable deformation without any improvement in the promotion metrics. The null result is therefore a regime-level null, not a measurement-layer absence.

\section{Persistence Metrics and Refinement Logic}

This section should give the paper a compact metric vocabulary that can be reused throughout. A useful drafting strategy is to present the persistence ladder in prose first, then insert a concise table of definitions and promotion criteria. The paper does not need to overformalize the metrics, but it should make clear that the same ladder is applied across all mechanism families.

The refinement logic should be described as part of the falsification discipline. If a branch shows bounded but non-promoting activity, it is recorded as active yet persistence-null. If a branch shows no meaningful activity at all, it is recorded separately as under-activated. Phase II is important precisely because many branches fall into the first category: they are active, bounded, and still negative. This distinction is what prevents the survey from collapsing into an implementation-null reading.

\begin{table}[h]
\centering
\small
\begin{tabular}{@{}p{2.0in}p{1.8in}p{2.2in}@{}}
\toprule
Metric & Operational role & Promotion use \\
\midrule
Composite lifetime & duration of local assembly & must increase before exploratory refinement \\
Binding persistence & retained close-configuration fraction & must increase for any persistence claim \\
Coarse basin persistence & mesoscopic residence stability & screens effective-medium style claims \\
Corridor dwell / locking & channel retention or transport anchoring & screens directional or gating claims \\
Stabilized ordering class & emergence of a new robust regime & required for regime-level reclassification \\
\bottomrule
\end{tabular}
\caption{Placeholder Table 1. Persistence metrics and promotion criteria. Final wording may be tightened, but the table should remain architecture-first and persistence-first.}
\label{tab:persistence_metrics}
\end{table}

\section{Mechanism Family Survey}

\subsection{Local Scalar Reinforcement (Stage 19A)}

This subsection should define Stage 19A as the minimal weak substrate-law test. The mechanism is a local scalar back-reaction in which kernel response is tied to local amplitude or energy-density-like occupation. The useful heuristic correspondence is to weak nonlinear, refractive, or effective-potential response, but the paper should explicitly say that this is a category correspondence only and not a physical identification \cite{nonlinear_placeholder}.

The observed response should be described as bounded and real. The local substrate field changes in a numerically controlled manner, and the branch should be presented as active even when the deformation remains small. A good placeholder sentence is: exact deformation norms and representative control-to-perturbed comparisons will be inserted from the Stage 19A run tables in the final draft.

The persistence result is nevertheless null. Composite lifetime, binding persistence, and coarse basin retention do not improve at the tested couplings. The subsection should therefore end with a restrained boundary statement: weak local scalar reinforcement, in the present architecture, is insufficient to produce persistence-level regime change.

\subsection{Mesoscopic Scalar Reinforcement (Stage 19B)}

This subsection should split the mesoscopic branch into basin-residence reinforcement and transport-corridor reinforcement while presenting them as one mechanism family. The heuristic category correspondence here is effective-medium or coarse geometric response: the law no longer follows only instantaneous local occupation, but instead responds to mesoscopic residence or channel usage over a structured region \cite{effective_medium_placeholder,emergent_geometry_placeholder}.

The bounded response is stronger here than in the purely local branch. Basin-residence reinforcement and corridor reinforcement both produce measurable mesoscopic fields and nontrivial substrate participation. That matters because it shows the architecture can register mesoscopic history-like structure without becoming trivial or numerically unstable. A placeholder sentence can note that the final version should insert representative field norms, basin dwell fractions, and kernel deviation summaries from the Stage 19B data products.

The persistence-level result remains null. Neither branch produces a gain in the primary observables, and no refinement promotion is earned. This is scientifically important because it weakens the case for ``more substrate response of the same basic kind'' as the next move. The null survives even when the substrate response is broadened from local to mesoscopic support.

\subsection{Directional Reinforcement (Stage 19C)}

Stage 19C changes symmetry class rather than simply scalar strength. The mechanism is directional or flux-driven substrate back-reaction. The cautious heuristic analogue is a gauge-like, Lorentz-like, or antisymmetric steering response, again framed only as a category correspondence and not as a claim of reproduced physics \cite{gauge_placeholder}.

The observed response should be described as bounded directional activity. The architecture registers directional support, curvature-proxy response, or transport-anisotropy signatures in a controlled way. This is enough to show that the directional branch is not implementation-null, even if the net deformation remains modest in the final tabulated values.

The persistence ladder does not move. No durable corridor locking, no anchoring, and no promoted stabilized transport class appear. The subsection should make the interpretive point clearly: directional response can be present without producing persistence. This distinguishes weak steering from genuine trapping or regime change.

\subsection{Mixed Scalar--Vector Reinforcement (Stage 19D)}

Stage 19D is the natural culmination of the Stage 19 family. It combines scalar-density and directional-flux reinforcement in one bounded law. The mechanism should be described as the strongest weak-coupling reinforcement attempt within Phase II, since it allows both occupancy-like and directional channels to contribute simultaneously.

The bounded response is again real. Mixed-field norms, kernel deviation, and transport-side diagnostics are non-null, and the final paper should insert representative values from the Stage 19D tables to show that the mixed branch is not simply a restatement of weaker pilots. This is the right place to emphasize that the negative survey is not built on inactive perturbations.

The persistence-level result, however, remains null. Even the strongest combined deformation in the Phase II reinforcement family fails to create persistent composites, anchoring, or promoted structure. This subsection should make explicit that the mixed branch closes the most obvious nearby extension of scalar and directional reinforcement taken together.

\subsection{Mutual Registration and Gated Interaction (Stage 21)}

Stage 21 leaves substrate reinforcement space and instead tests selective interaction eligibility. The mechanism family includes phase-coherence registration, spatial-configuration registration, and composite-topology registration. The heuristic analogue is a selection-rule, resonance-window, or interaction-eligibility law: interaction is not always ``on,'' but depends on a compatibility measure \cite{selection_rule_placeholder,resonance_placeholder}.

The observed response is one of the strongest examples of active-yet-negative behavior. Eligibility windows can open, duty cycles can become nonzero or even sustained, and effective interaction fractions can be measured. Spatial and topological registration classes are especially useful here because they show that the architecture can recognize compatibility structure without turning it into persistence.

The persistence result remains null across the registration classes. No branch produces a gain in composite lifetime, binding persistence, basin retention, or corridor locking. The structural point is that compatibility detection alone is not enough. The architecture can implement a gated interaction rule without generating a new persistence regime.

\subsection{Delayed Nonlocal Memory (Stage 22)}

Stage 22 tests whether temporally accumulated nonlocal memory can supply what instantaneous weak reinforcement and instantaneous registration do not. The mechanism family includes single-lag retarded response, exponential-history accumulation, and finite-window averaging. The cautious category correspondence is non-Markovian linear response or retarded effective-medium behavior \cite{nonmarkov_placeholder,retarded_response_placeholder}.

The observed response is again bounded and non-null. Memory fields are activated, historical accumulation becomes measurable, and delayed substrate deviation remains numerically controlled. In the final draft this subsection should insert the representative comparison that shows the exponential-history class is stronger than the single-lag and finite-window classes while remaining persistence-null.

The persistence-level result still does not change. No anchoring, no retained composite, and no promoted stabilized class appear. This is one of the sharpest negative results in the Phase II survey because it closes a natural class of ``history should help'' intuitions without requiring overclaiming about all possible memory laws.

\section{Relation to Established Physics Categories}

This section should explicitly state that the mechanism families are being related to known physics-style categories only heuristically. Local scalar reinforcement is category-adjacent to weak nonlinear or effective-potential response; mesoscopic reinforcement is category-adjacent to effective-medium or emergent geometric response; directional reinforcement is category-adjacent to gauge-like steering analogues; registration gating is category-adjacent to resonance-window or selection-rule logic; delayed memory is category-adjacent to non-Markovian or retarded-response models.

The paper should also say what is \emph{not} being claimed. These correspondences are not proofs of equivalence, not demonstrations that the architecture reproduces known physics, and not evidence that any real-world interaction has been derived. They function only as a way of locating the tested law classes within a familiar conceptual map. That restrained positioning is important for preventing analogy from doing argumentative work it cannot support.

\begin{table}[h]
\centering
\small
\begin{tabular}{@{}p{1.6in}p{2.0in}p{2.2in}@{}}
\toprule
Mechanism family & Heuristic analogue category & Interpretive restriction \\
\midrule
Local scalar reinforcement & weak nonlinear / effective-potential response & category correspondence only \\
Mesoscopic reinforcement & effective medium / coarse geometric response & not an emergent geometry claim \\
Directional reinforcement & gauge-like or Lorentz-like steering analogue & not a gauge theory derivation \\
Registration gating & resonance-window / selection-rule eligibility & not a quantum selection rule claim \\
Delayed memory & non-Markovian / retarded-response analogue & not a field-theoretic memory derivation \\
\bottomrule
\end{tabular}
\caption{Placeholder Table 3. Mapping from tested mechanism classes to cautious physics-style analogue categories.}
\label{tab:physics_mapping}
\end{table}

\section{Structural Boundary Statement}

This section is the heart of the paper and should be written as plainly as possible. A good central statement is:
\begin{quote}
Across the tested coupling scales and symmetry classes, weak substrate-law and gated-interaction modifications do not generate persistent composites or anchoring in the present kernel-induced cochain architecture.
\end{quote}

The surrounding prose should clarify why this is not a trivial null. Many branches produce bounded non-null deformation, bounded interaction eligibility, or bounded temporal accumulation. The architecture responds. What remains absent is the regime-level change in the persistence observables. The null is therefore structural rather than microscopic inactivity.

This is also the right place to incorporate the role of Stage 20. Rather than standing alone as a numerical section, Stage 20 functions here as the conceptual consolidation step. It marks the transition from ``several negative pilots'' to ``a mapped no-law boundary'' for a broad nearby family of weak classical or bosonic-style response mechanisms.

\begin{table}[h]
\centering
\small
\begin{tabular}{@{}p{1.0in}p{2.0in}p{2.6in}@{}}
\toprule
Stage & Mechanism class & Consolidated null result \\
\midrule
19A & local scalar back-reaction & bounded local substrate response, no persistence gain \\
19B & mesoscopic reinforcement & active basin/corridor response, no anchoring or promotion \\
19C & directional reinforcement & bounded steering-like response, no locking regime \\
19D & mixed scalar-vector reinforcement & strongest combined weak deformation, still persistence-null \\
21 & registration-gated interaction & active eligibility windows, no persistence-level change \\
22 & delayed nonlocal memory & active temporal accumulation, no retained composite \\
\bottomrule
\end{tabular}
\caption{Placeholder Table 4. Consolidated null-result summary across the completed Phase II survey.}
\label{tab:null_summary}
\end{table}

\section{Discussion}

The discussion should state carefully what has and has not been falsified. The survey does not falsify all possible nonlinear laws, all possible operator classes, or all possible paths to composite persistence. It does not demonstrate that no future extension can work. What it does rule out is a broad nearby family of weak substrate-law, weak gating, and weak delayed-memory mechanisms within the present architectural baseline.

This distinction matters for scientific pacing. A negative survey is useful only if it shrinks the nearby search space. In the present case it suggests that additional coefficient chasing within the same weak-law family is unlikely to be the most informative next move. The architecture has already shown that bounded response can exist without persistence. More of the same response class is therefore not, by itself, a compelling next hypothesis.

The discussion should then motivate future directions cautiously. The most natural next candidates are those that change what kind of object is propagating or how the operator sector is organized: new degrees of freedom, multi-component or graded states, stronger constraint classes, topological structure, or mediator-like relational channels. These are motivations only. They are not claims that such extensions have succeeded or must succeed.

\section{Conclusion}

The conclusion should be short and explicit. Phase II remains persistence-null across the tested weak mechanism family. The survey nevertheless matters because the negative is disciplined, cross-checked, and architecture-level. Weak substrate reinforcement, selective interaction eligibility, and delayed nonlocal memory all enter the model in bounded and numerically real ways without producing persistent composites or anchoring.

The resulting interpretation is that Phase II defines a stable no-law boundary for the present architectural neighborhood. That boundary, rather than a single null stage, is the paper's main contribution. It motivates a transition away from nearby weak-law variants and toward later Phase III sector change.

\section{Outlook}

The outlook section should remain brief and noncommittal. Natural continuations include the Dirac--K\"ahler or graded operator sector, multi-component packet states, topological or exclusion-like constraint structures, and mediator or relational-state channels \cite{dirac_kahler_placeholder}. The point is not that these mechanisms are already validated, but that the completed negative survey justifies looking beyond the weak classical / bosonic-style law family tested here.

Future work should therefore be framed as a change in sector, degree of freedom, or constraint class rather than another nearby reinforcement pilot. That is the cleanest way to respect the evidence accumulated by the Phase II boundary program.

\section{Placeholder Figures and Tables}

The final paper will likely benefit from a small number of summary figures and tables rather than a large figure stack. The placeholders below are intended to reserve space for the core visual argument.

\begin{figure}[h]
\centering
\placeholderbox{1.8in}
\caption{Placeholder Figure 1. Schematic of the kernel-induced cochain architecture and the persistence-first observable layer. This figure should explain the fixed baseline, the representative packet families, and the promotion ladder.}
\label{fig:architecture}
\end{figure}

\begin{figure}[h]
\centering
\placeholderbox{1.8in}
\caption{Placeholder Figure 2. Summary matrix of tested mechanism families and persistence outcomes across Stages 19A, 19B, 19C, 19D, 21, and 22.}
\label{fig:summary_matrix}
\end{figure}

\begin{figure}[h]
\centering
\placeholderbox{1.8in}
\caption{Placeholder Figure 3. Example bounded-but-persistence-null deformation plots, showing that numerically real substrate response does not imply regime change.}
\label{fig:bounded_null}
\end{figure}

\begin{figure}[h]
\centering
\placeholderbox{1.8in}
\caption{Placeholder Figure 4. Delayed-memory class comparison for single-lag, exponential-history, and finite-window kernels.}
\label{fig:memory_comparison}
\end{figure}

\begin{figure}[h]
\centering
\placeholderbox{1.8in}
\caption{Placeholder Figure 5. Registration-class activity versus null persistence outcome, contrasting active eligibility windows with unchanged primary observables.}
\label{fig:registration_null}
\end{figure}

\begin{table}[h]
\centering
\small
\begin{tabular}{@{}p{1.2in}p{1.8in}p{2.6in}@{}}
\toprule
Table & Purpose & Notes \\
\midrule
Table 1 & persistence metrics and promotion criteria & see Table~\ref{tab:persistence_metrics} \\
Table 2 & mapping from stages to mechanism classes & insert finalized stage-class summary here \\
Table 3 & mechanism classes to analogue categories & see Table~\ref{tab:physics_mapping} \\
Table 4 & consolidated null-result summary & see Table~\ref{tab:null_summary} \\
\bottomrule
\end{tabular}
\caption{Placeholder Table list for the final paper assembly pass.}
\end{table}

\appendix

\section{Appendix Placeholder: Representative Run-Sheet Conventions}

This appendix can collect representative run-sheet fields, branch identifiers, and control-versus-perturbed naming conventions so that the main text stays readable while the numerical workflow remains auditable.

\section{Appendix Placeholder: Numerical Stability and Resolution Policy}

This appendix can summarize timestep policy, bounded-coupling policy, control matching, and the criteria used to distinguish active bounded response from under-activated branches.

\section{Appendix Placeholder: Stage-to-Stage Metric Definitions}

This appendix can give compact operational definitions for the persistence ladder and document any stage-local derived quantities that appear in summary tables.

\begin{thebibliography}{99}

\bibitem{dec_placeholder}
Placeholder reference on discrete exterior calculus, lattice gauge structure, or related cochain operator formalisms.

\bibitem{nonlinear_placeholder}
Placeholder reference on nonlinear wave dynamics, weak refractive response, or effective-potential analogues.

\bibitem{effective_medium_placeholder}
Placeholder reference on effective media, homogenization, or coarse-grained response in structured transport systems.

\bibitem{emergent_geometry_placeholder}
Placeholder reference on emergent metric or geometry-like response in effective discrete systems.

\bibitem{gauge_placeholder}
Placeholder reference on gauge-like transport, antisymmetric steering response, or lattice transport analogues.

\bibitem{selection_rule_placeholder}
Placeholder reference on selection rules, gated interactions, or compatibility-controlled coupling.

\bibitem{resonance_placeholder}
Placeholder reference on resonance windows, transport locking, or interaction eligibility in coupled wave systems.

\bibitem{nonmarkov_placeholder}
Placeholder reference on non-Markovian field theory, memory kernels, or delayed-response dynamics.

\bibitem{retarded_response_placeholder}
Placeholder reference on retarded effective media or temporally accumulated linear response.

\bibitem{dirac_kahler_placeholder}
Placeholder reference on Dirac--K\"ahler operators, graded packet states, or discrete fermionic analogues.

\end{thebibliography}

\end{document}
